% !TEX root = ./main.tex

\chapter*{Introduction}
\addcontentsline{toc}{chapter}{Introduction}

Depuis la fin des années 1990, les dérivés de crédit occupent une place de
plus en plus importante dans la gestion et le transfert du risque de crédit.
Parmi ces instruments, le \gls{cds} s'est imposé comme
le produit de référence, permettant à un acheteur de protection de se
prémunir contre le risque de défaut d'une entité tierce moyennant le
versement d'une prime périodique à un vendeur de protection. Au-delà de
son rôle de couverture, le \gls{cds} constitue également un indicateur de marché
précieux : le spread qui lui est associé reflète en temps réel la
perception, par les investisseurs, de la qualité de crédit de l'entité de
référence.

La valorisation de ces contrats repose nécessairement sur une modélisation
mathématique du risque de défaut. Deux grandes familles d'approches
coexistent dans la littérature : les modèles structurels, issus des travaux
de Merton, qui relient le défaut à la valeur des actifs de l'entreprise, et
les modèles à intensité, dans lesquels le défaut est modélisé comme l'instant de saut d'un processus de comptage dont
l'intensité, appelée intensité de défaut, est elle-même un processus
stochastique. Cette seconde approche présente l'avantage de ne pas exiger
d'hypothèses sur la structure du bilan de l'entreprise et de se prêter
naturellement à la valorisation des produits dérivés de crédit à partir
des données de marché observables.

Dans ce cadre, le choix du processus stochastique retenu pour représenter
l'intensité de défaut est déterminant. Ce processus doit en effet respecter
certaines propriétés économiques essentielles : rester positif à tout
instant, puisqu'il s'agit d'un taux instantané de défaut, et présenter un
comportement de retour vers un niveau d'équilibre de long terme, conformément
aux observations empiriques sur l'évolution du risque de crédit. Le modèle
de \gls{cir}, initialement développé pour la modélisation
des taux d'intérêt courts, satisfait ces deux exigences et offre de surcroît
une tractabilité analytique remarquable grâce à sa structure affine, ce qui
en fait un candidat naturel pour la modélisation de l'intensité de défaut.

L'objectif de ce mémoire est double. Il s'agit, d'une part, de construire un
cadre probabiliste rigoureux permettant de modéliser le défaut d'une entité
de référence à l'aide d'un processus d'intensité de type \gls{cir}, puis d'en
déduire les formules de valorisation des deux jambes d'un \gls{cds}, la jambe fixe
et la jambe de protection, ainsi que le spread théorique du contrat.
D'autre part, il s'agit de proposer et de mettre en \oe uvre une méthode
d'estimation des paramètres du modèle à partir des spreads de \gls{cds}
observés sur le marché, puis d'appliquer cette méthode à des données
réelles afin d'évaluer la capacité du modèle à reproduire les
observations de marché et d'en discuter les forces et les limites.

Pour répondre à ces objectifs, ce mémoire est organisé en trois parties.
La première partie pose les bases théoriques nécessaires à l'étude :
elle présente le fonctionnement des Credit Default Swaps, rappelle les
notions fondamentales de théorie des probabilités et de calcul stochastique
utiles à la suite du mémoire, et introduit le modèle \gls{cir} ainsi que les
principales méthodes d'estimation statistique. La deuxième partie
développe le cadre de modélisation du défaut dans un espace probabilisé
filtré, construit le modèle \gls{cir} multifactoriel retenu pour l'intensité de
défaut, établit les formules de valorisation des deux jambes du \gls{cds}, puis
propose une méthode d'estimation des paramètres du modèle fondée sur la
méthode des moments appliquée à la distribution stationnaire du processus
\gls{cir}. Enfin, la troisième partie est consacrée à l'application numérique du
modèle sur des données réelles de spreads de \gls{cds} pour trois entreprises
présentant des profils de risque de crédit contrastés, à l'analyse des
résultats obtenus, à une étude de sensibilité des paramètres, ainsi
qu'à une discussion critique des forces et des limites du modèle \gls{cir}.